% !PRINTED 403
\weberpage{403}

Hiernach ergibt sich aus (2):

\weq{5}{\begin{aligned}
\mathrm{Lim}_{t=\infty}\frac{Z_t}{t} &= \frac{\log\varepsilon}{a\sqrt D}, & D&>0,\\
&= \frac{2\pi}{a\sqrt{-D}}, & D&<0.
\end{aligned}}
Nach Bd. II, § 196, 4. ist aber

\[ \mathrm{Lim}_{s=1}\sum\frac{s-1}{(Ny)^s} = \mathrm{Lim}_{t=\infty}\frac{Z_t}{t}, \]
und es ergibt sich also aus (5):

\weq{6}{\begin{aligned}
\mathrm{Lim}_{s=1}\sum\frac{s-1}{(Ny)^s} &= \frac{1}{a\sqrt D}\log\varepsilon, & D&>0,\\
&= \frac{2\pi}{a\sqrt{-D}}, & D&<0.
\end{aligned}}
Hier durchläuft $y$ die Reihe der ganzen Zahlen der Ordnung
$[Q]$, die durch ein zu $Q$ teilerfremdes primäres Ideal $\mathfrak{a}$ teilbar
sind, wobei jedoch von einer Schar assoziierter Zahlen immer
nur ein Repräsentant, der jetzt beliebig gewählt sein kann, beizubehalten
 ist (bei negativen $D$ zwei oder vier oder sechs).

Hierbei sind unter assoziierten Zahlen (nach der Ordnung
$[Q]$) solche zu verstehen, deren Quotient \gesperrt{eine Einheit der
Ordnung $Q$ ist.}

Wir wollen jetzt von den Zahlen $y$ noch alle die ausschließen,
die nicht relativ prim zu $Q$ sind. Dann modifiziert sich die
Summe (6) in folgender Weise: Ist $r$ ein Primfaktor von $Q$, so
sind alle Zahlen der Ordnung $[Q]$, die zu $r$ nicht teilerfremd
sind, durch $r$ teilbar, weil sie ja nach dem Modul $r$ mit einer
rationalen Zahl kongruent sind.

Setzen wir also

\weq{7}{y = r\,y_1,}
so ist $y_1$ durch $\mathfrak{a}$ teilbar, braucht aber nicht der Ordnung $[Q]$,
sondern nur der Ordnung $\left[\dfrac{Q}{r}\right] = [Q_1]$ anzugehören. Ist bei
positiver Diskriminante $\varepsilon_1$ die fundamentale Einheit der Ordnung
$[Q_1]$, so ist

\weq{8}{\varepsilon = \varepsilon_1^{\lambda}}
eine Potenz von $\varepsilon_1$, und es ist $\sqrt D = r\sqrt{D_1}$. Die Zahlen

\weq{9}{y = r\,y_1\,\varepsilon_1^{t}, \quad t = 0, 1, 2, \ldots, \lambda - 1}
%CONT>>
